\section{\texttt{Data.Real.Algebra.Addition}}

\begin{lemma}[\agda{+InvR}]
  \label{+InvR}
  For every $x : \R$, $x + (-x) = 0$.
\end{lemma}
\begin{proof}
  By \Cref{composeNE2} applied to the identity, negation and addition, the map $x \mapsto x + (-x)$ is $2$-Lipschitz. The constant map at $0$ is also Lipschitz, so \Cref{lipschitzEq} reduces the claim to rational points, where the computation rules in \Cref{liftNE,liftNE2} and the inverse law in $\Q$ give
  \[
    \iota(q) + (-\iota(q)) = \iota(q + (-q)) = \iota(0).
  \]
\end{proof}

\begin{corollary}[\agda{ℝAbGroup}; {\textcite[\S~11.3.3]{univalentfoundationsprogram2013}}]
  \label{RAbGroup}
  $(\R, 0, +, -)$ is an abelian group.
\end{corollary}
\begin{proof}
  $\R$ is a set, and addition is associative and commutative and satisfies the left and right unit laws with respect to $0$, each established alongside the definition of the operations by \Cref{nonExpansive2Eq}. Together with \Cref{+InvR} these are the abelian group axioms.
\end{proof}
